\documentclass[12pt,a4paper]{article}

\usepackage[T1]{fontenc}
\usepackage[utf8]{inputenc}
\usepackage[lithuanian]{babel}
\usepackage{amsmath}
\usepackage{booktabs}
\usepackage{parskip}
\usepackage{geometry}
\geometry{top=2.5cm,bottom=2.5cm,left=3cm,right=2.5cm}

\begin{document}

% ===== TITULINIS =====

\begin{center}
    {\large\textbf{Vilniaus Universitetas}}\\[0.2cm]
    {\large Matematikos ir Informatikos Fakultetas}\\[1.5cm]
    {\LARGE\bfseries Neuroninis tinklas - Iris gėlės klasifikavimas}\\[0.3cm]
    {\large Laboratorinio darbo ataskaita}\\[2cm]

    \begin{tabular}{rl}
        \textbf{Studentas:}          & Pijus Ranonis \\[0.2cm]
        \textbf{Kursas / grupė:}     & 1 kursas, 5 grupė \\[0.2cm]
        \textbf{Pratybų dėstytojas:} & Linas Litvinas \\
        \textbf{Dalyko dėstytojas:}  & Rimantas Vaicekauskas \\
    \end{tabular}

    \vspace{2cm}
    {\large Vilnius, \the\year}
\end{center}
\newpage

% ===== APRAŠYMAS ======

\section*{Darbo tikslas ir aprašymas}

Šio laboratorinio darbo tikslas - sukurti programą, kuri pasinaudojus dirbtinį neuroninį tinklą identifikuoja Iris gėlės rūšį.

Gėlės rūšys:
\begin{itemize}
    \item \textit{Iris setosa} (reikšmė: < 1,5)
    \item \textit{Iris versicolor} (reikšmė: 1,5 -- 2,5)
    \item \textit{Iris virginica} (reikšmė: > 2,5)
\end{itemize}

\section*{Architektūra ir veikimo etapai}

\subsection*{1 etapas: Tinklo sudarymas}

Tinklas sudarytas iš dviejų sluoksnių:

\begin{table}[h]
\centering
\caption{Neuroninio tinklo sluoksnių parametrai}
\begin{tabular}{@{}llll@{}}
\toprule
\textbf{Sluoksnis} & \textbf{Įėjimai} & \textbf{Išėjimai} & \textbf{Aktyvacijos funkcija} \\
\midrule
Įėjimo sluoksnis    & 4 & 6 & Sigmoid \\
Paslėptas sluoksnis & 6 & 1 & Tiesinė \\
\bottomrule
\end{tabular}
\end{table}

\subsection*{2 etapas: Duomenų normalizavimas}

Prieš mokymą visi matavimų duomenys normalizuojami į intervalą $[0,\;1]$ naudojant formulę:
\[
x' = \frac{x - x_{\min}}{x_{\max} - x_{\min}}
\]

\subsection*{3 etapas: Evoliucinis mokymo algoritmas}

Tinklas treniruojamas evoliucijos algoritmo:
\begin{enumerate}
    \item Sukuriamos atsitiktinės svorių reikšmė.
    \item Kiekvienoje iteracijoje generuojami atsitiktiniai svorių ir poslinkių pokyčiai.
    \item Apskaičiuojamas naujo tinklo MSE su visais 150 pavyzdžių.
    \item Jei naujasis MSE mažesnis -- naujas tinklas perrašomas vietoj senojo.
    \item Pakeitimų dydis $\Delta$ palaipsniui mažinamas su kiekviena iteracija (nuo 0,2 iki 0).
\end{enumerate}

\subsection*{Klaidos funkcija}

Tinklo kokybei vertinti naudojamas vidutinis kvadratinės paklaidos rodiklis (MSE):
\[
\text{MSE} = \frac{1}{n}\sum_{i=1}^{n}\left(\hat{y}_i - y_i\right)^2
\]
kur $\hat{y}_i$ -- tinklo prognozuojama reikšmė, $y_i$ -- tikroji reikšmė, $n$ -- pavyzdžių skaičius.

\section*{Gauti rezultatai}

\begin{table}[h]
\centering
\caption{MSE reikšmės prieš ir po treniravimo}
\begin{tabular}{@{}lcc@{}}
\toprule
\textbf{Etapas} & \textbf{MSE} & \textbf{Paklaidos sumažėjimas} \\
\midrule
Pradinis tinklas (atsitiktiniai svoriai) & $\approx 20{,}90$ & -- \\
Po treniravimo (100\,000 iteracijų)      & $\approx 0{,}02$  & $\approx 99{,}9\,\%$ \\
\bottomrule
\end{tabular}
\end{table}

Treniravimo metu MSE sumažėjo nuo $\approx 20{,}90$ iki $\approx 0{,}02$ -- beveik 1000 kartų mažiau.

\section*{Išvados}

\begin{enumerate}
    \item Sėkmingai suprogramuotas dviejų sluoksnių dirbtinis neuroninis tinklas Java kalba.
    \item Evoliucinio mokymo algoritmas veikia ir gali pasiekti žemą MSE ($\approx 0{,}02$) per pakankamai iteracijų.
\end{enumerate}

\end{document}